\documentclass[12pt,a4paper]{article}

% ---- Encoding & Font ----
\usepackage{iftex}
\ifPDFTeX
  \usepackage[utf8]{inputenc}
  \usepackage[T1]{fontenc}
\else
  \usepackage{fontspec}  % XeTeX/LuaTeX (tectonic): Unicode nativo (·, —, ¿, ª, §)
\fi
\usepackage[spanish]{babel}
\usepackage{csquotes}

% ---- Page layout ----
\usepackage[top=2.5cm, bottom=2.5cm, left=2.5cm, right=2.5cm]{geometry}
\usepackage{setspace}
\onehalfspacing

% ---- Math ----
\usepackage{amsmath, amssymb, amsthm}
\usepackage{mathtools}
\usepackage{physics}
\usepackage{esint}  % \oiint

% ---- Graphics ----
\usepackage{graphicx}
\usepackage{tikz}
\usepackage{float}
\usepackage{caption}

% ---- Tables ----
\usepackage{array}
\usepackage{booktabs}
\usepackage{tabularx}
\usepackage{colortbl}

% ---- Colours ----
\usepackage{xcolor}
\definecolor{notebg}{HTML}{FFF3CD}
\definecolor{noteborder}{HTML}{FFC107}
\definecolor{boxred}{HTML}{F8D7DA}
\definecolor{boxredborder}{HTML}{F5C6CB}

% ---- Boxes ----
\usepackage{mdframed}
\usepackage{tcolorbox}
\tcbuselibrary{most}

\newtcolorbox{keybox}[1][]{
  colback=blue!5,
  colframe=blue!70!black,
  arc=4pt,
  boxrule=1pt,
  fonttitle=\bfseries,
  title={#1}
}

\newtcolorbox{warnbox}[1][]{
  colback=boxred,
  colframe=boxredborder,
  arc=4pt,
  boxrule=1pt,
  fonttitle=\bfseries,
  title={#1}
}

\newtcolorbox{notebox}[1][]{
  colback=notebg,
  colframe=noteborder,
  arc=4pt,
  boxrule=1pt,
  fonttitle=\bfseries,
  title={#1}
}

% ---- Hyperlinks & TOC ----
\usepackage[hidelinks]{hyperref}
\usepackage{bookmark}

% ---- Enumerate ----
\usepackage{enumitem}

% ---- Miscellaneous ----
\usepackage{icomma}
\usepackage{siunitx}
\usepackage{textcomp}
\usepackage[bottom]{footmisc}

% ---- Title ----
\title{\textbf{Clase 16: Par sobre una Espira, Motores y Galvanómetro, \\
        y el Campo Magnético de las Corrientes (Biot--Savart)}\\[0.3em]
        \large{Torque sobre una espira, instrumentos de medida y campo
              generado por cargas y corrientes}}
\author{Transcripto y expandido --- Curso Física III (OpenFING)}
\date{}

\begin{document}

\maketitle
\thispagestyle{empty}
\newpage

\tableofcontents
\newpage

% ===================================================================
\section{Par sobre una espira con corriente}
% ===================================================================

\subsection{Configuración}

Espira rectangular con corriente $I$ en un campo $\mathbf{B}$
\textbf{uniforme}, que gira sobre un eje. Versor normal $\hat{\mathbf{n}}$
y ángulo $\theta$ con $\mathbf{B}$.

\begin{notebox}[Campo no uniforme]
  Si el campo no es uniforme pero la espira es muy pequeña, se lo trata
  como uniforme.
\end{notebox}

\subsection{Análisis de las fuerzas}

Fuerza sobre cada lado $\mathbf{F} = I\,\mathbf{L}\times\mathbf{B}$:
\begin{itemize}[nosep]
  \item \textbf{Lados horizontales:} fuerzas verticales opuestas, se
        anulan; líneas de acción sobre el eje, no hacen par.
  \item \textbf{Lados laterales:} $\mathbf{F}_3, \mathbf{F}_4$ opuestas
        (resultante nula) pero a distinto lado del eje: generan
        \textbf{par}.
\end{itemize}
Resultante cero, pero \textbf{torque} neto. (Con campo no uniforme
habría también resultante.)

\subsection{El torque y su módulo}

Lados $a$, $b$ ($A=ab$). Torque de $\mathbf{F}_3$ (brazo
$\tfrac{a}{2}\sin\theta$), con $F_3 = IbB$; sumando $\mathbf{F}_4$:

\begin{keybox}[Torque sobre una espira]
  \[
  \boxed{\tau = I\,A\,B\,\sin\theta}
  \]
  Nulo en $\theta=0$ (normal $\parallel$ campo): equilibrio alineado.
\end{keybox}

\subsection{Forma vectorial: vector área y momento}

Vector área $\mathbf{A} = A\,\hat{\mathbf{n}}'$ con $\hat{\mathbf{n}}'$
por la regla de la mano derecha (dedos $=$ corriente, pulgar $=$ normal):

\begin{keybox}[Torque vectorial]
  \[
  \boxed{\boldsymbol{\tau} = I\,\mathbf{A}\times\mathbf{B}}
  \]
  (orden $\mathbf{A}\times\mathbf{B}$ para el sentido correcto).
\end{keybox}

\begin{notebox}[Momento dipolar magnético]
  $\boldsymbol{\mu} = I\,\mathbf{A}$, así $\boldsymbol{\tau} =
  \boldsymbol{\mu}\times\mathbf{B}$, análogo a $\boldsymbol{\tau} =
  \mathbf{p}\times\mathbf{E}$.
\end{notebox}

\subsection{Bobina de N vueltas}

\[
\boxed{\boldsymbol{\tau} = N\,I\,\mathbf{A}\times\mathbf{B}}
\]

\begin{notebox}[Consejo]
  Calcular el módulo con la fórmula y el sentido pensando las fuerzas.
\end{notebox}

% ===================================================================
\section{Aplicaciones}
% ===================================================================

\subsection{Motor de corriente continua}

Sin más, la espira oscila (el torque se invierte al pasar la alineación).
Un \textbf{conmutador} (escobillas) invierte la corriente cada media
vuelta, manteniendo el giro: \textbf{motor de corriente continua}. El de
corriente alterna es igual, pero la corriente ya oscila sola.

\begin{notebox}[Torque variable]
  Como $\tau\propto\sin\theta$ no es constante, se agregan dispositivos
  estabilizadores.
\end{notebox}

\subsection{Galvanómetro: amperímetro y voltímetro}

Como $\tau\propto I$, con un \textbf{resorte de torsión} (torque
$\propto$ ángulo) se mide la corriente en la posición de equilibrio:
\textbf{galvanómetro}.

\begin{itemize}[nosep]
  \item \textbf{Amperímetro:} en \textbf{serie}, resistencia
        \textbf{pequeña} (no alterar la corriente).
  \item \textbf{Voltímetro:} en \textbf{paralelo}, resistencia
        \textbf{grande} (desviar poca corriente); $\Delta V = RI$.
\end{itemize}

\begin{keybox}[Principio de medición]
  Todo instrumento perturba; se lo diseña para perturbar lo menos
  posible.
\end{keybox}

% ===================================================================
\section{Campo magnético de una carga en movimiento}
% ===================================================================

Capítulo \textbf{9}: cómo las cargas en movimiento \textbf{generan}
campo.

\subsection{Dificultad experimental}

Una carga de prueba sentiría fuerza eléctrica \textbf{y} magnética, pero
la magnética es muchísimo menor: queda oculta. En la práctica se mide el
campo de una \textbf{corriente} (cable neutro). La carga puntual se
estudia por ser más fácil de describir.

\subsection{Descripción del campo}

En un círculo en el plano $\perp\mathbf{v}$ que pasa por $q$: $|\mathbf{B}|$
constante, tangente al círculo (mano derecha). Además $B\propto 1/r^2$,
$\propto v$, $\propto |q|$, $\propto \sin\theta$.

\subsection{Ley de Biot--Savart para una carga puntual}

\begin{keybox}[Campo de una carga en movimiento]
  \[
  \boxed{\mathbf{B}(\mathbf{r}) = \frac{\mu_0}{4\pi}\,
         \frac{q\,\mathbf{v}\times\mathbf{r}}{r^3}}
  \]
  $\mu_0$: permeabilidad magnética del vacío (análoga a $\varepsilon_0$).
\end{keybox}

\subsection{Comparación con Coulomb y validez}

Se parece a Coulomb ($\propto 1/r^2$, $\propto q$), pero con
\textbf{estructura vectorial} ($\perp\mathbf{r}$) y \textbf{dependencia de
la velocidad}.

\begin{warnbox}[Validez $v \ll c$]
  La fórmula es instantánea; en realidad las señales viajan a $c$. Si $v
  \sim c$, deja de valer.
\end{warnbox}

% ===================================================================
\section{Ley de Biot--Savart para un cable}
% ===================================================================

\subsection{Deducción}

Elemento $d\mathbf{s}$ con corriente $I$; con $dQ\,\mathbf{v}_d = I\,d\mathbf{s}$
(pues $d\mathbf{s} = \mathbf{v}_d\,dt$, $I = dQ/dt$):
\[
d\mathbf{B} = \frac{\mu_0}{4\pi}\,\frac{I\,d\mathbf{s}\times\mathbf{r}}{r^3}.
\]

\begin{keybox}[Ley de Biot--Savart]
  \[
  \boxed{\mathbf{B} = \frac{\mu_0}{4\pi}\,I\int_C
         \frac{d\mathbf{s}\times\mathbf{r}}{r^3}}
  \]
  Da el campo conocidas las fuentes (corrientes).
\end{keybox}

\subsection{Ejemplo: segmento rectilíneo}

Cable de largo $L$, corriente $I$; campo en $P$ a distancia $X$ en el
plano bisector. Todas las contribuciones $d\mathbf{B}$ son entrantes y
\textbf{paralelas} $\Rightarrow$ se suman los módulos. Con $\sin\theta =
X/\sqrt{X^2+Y^2}$, $r^2 = X^2+Y^2$:
\[
dB = \frac{\mu_0 I}{4\pi}\,\frac{X\,dy}{(X^2+Y^2)^{3/2}}
\;\Longrightarrow\;
B = \frac{\mu_0 I X}{4\pi}\int_{-L/2}^{L/2}\frac{dy}{(X^2+Y^2)^{3/2}}.
\]

\begin{notebox}[Cálculo]
  Mismo tipo de integral que la barra cargada (Clase 8). Se resuelve con
  $Y = Xu$ y un truco de integración por partes (sumar y restar
  $1 = (1+u^2) - u^2$).
\end{notebox}

\subsection{Límites: dipolo magnético e hilo infinito}

\textbf{Lejos ($X \gg L$):} el segmento da $B\sim 1/X^2$, pero el circuito
cerrado completo decrece como $1/r^3$.

\begin{keybox}[Dipolo magnético]
  Sin monopolos magnéticos, un circuito visto de lejos es un
  \textbf{dipolo}: $B \sim 1/r^3$ (como el dipolo eléctrico, no como la
  carga puntual $1/r^2$).
\end{keybox}

\textbf{Cerca ($X \ll L$):} campo del hilo infinito:
\[
\boxed{B = \frac{\mu_0 I}{2\pi X}}.
\]

\begin{notebox}[Dirección]
  Líneas de $\mathbf{B}$ circulares cerradas alrededor del cable,
  tangentes, orientadas por la mano derecha (pulgar $=$ corriente).
\end{notebox}

\begin{notebox}[Próxima clase]
  Hilo infinito, fuerza entre corrientes y ley de Ampère.
\end{notebox}

\end{document}
